% Part: many-valued-logic
% Chapter: syntax-and-semantics
% Section: connectives

\documentclass[../../../include/open-logic-section]{subfiles}

\begin{document}

\olfileid{mvl}{syn}{con}

\olsection{اللغات والروابط}

لمنطق القضايا الكلاسيكي، ولكثير من المنطقات غيره، ذخيرة معينة من
\emph{الثوابت القضوية} و\emph{الروابط}. فمن الرموز الأولية التي
نستعملها:
\begin{enumerate}
\tagitem{prvFalse}{الثابت القضوي الدال على~!!{falsity}، وهو~$\lfalse$.}{}
\tagitem{prvTrue}{الثابت القضوي الدال على~!!{truth}، وهو~$\ltrue$.}{}
\item الروابط المنطقية:
  \startycommalist
  \iftag{prvNot}{\ycomma $\lnot$ (النفي)}{}%
  \iftag{prvAnd}{\ycomma $\land$ (الاقتران)}{}%
  \iftag{prvOr}{\ycomma $\lor$ (الفصل)}{}%
  \iftag{prvIf}{\ycomma $\lif$ (!!{conditional})}{}%
  \iftag{prvIff}{\ycomma $\liff$ (!!{biconditional})}{}%
\end{enumerate}
\iftag{defNot,defOr,defAnd,defIf,defIff,defTrue,defFalse,defEx,defAll}{%
وفوق هذه الروابط الأولية قد نستعمل رموزًا معرفة بطريق الاختصار، ومنها
\startycommalist
  \iftag{defNot}{\ycomma $\lnot$ (النفي)}{}%
  \iftag{defAnd}{\ycomma $\land$ (الاقتران)}{}%
  \iftag{defOr}{\ycomma $\lor$ (الفصل)}{}%
  \iftag{defIf}{\ycomma $\lif$ (!!{conditional})}{}%
  \iftag{defIff}{\ycomma $\liff$ (!!{biconditional})}{}%
  \iftag{defFalse}{\ycomma $\lfalse$ (!!{falsity})}{}%
  \iftag{defTrue}{\ycomma $\ltrue$ (!!{truth})}.}{}

وهذه الروابط بعينها صالحة للمنطقات المتعددة القيم. إلا أننا قد نحتاج
في المنطق الواحد إلى صور متعددة من رابط واحد، كالاقتران، فلا بد من
رموز مختلفة لتمييزها. وقد تدخل في بعض المنطقات المتعددة القيم روابط
لا نظير لها في المنطق الكلاسيكي. فلذلك نأخذ اللغة هنا على وجه أعم.

\begin{defn}
  \emph{لغة القضايا} مجموعة~$\Lang L$ من \emph{الروابط}. ولكل رابط
  $\star$ \emph{رتبة}؛ فإن كانت رتبته~$\OLId{latin-ordinary}{n}$ قيل إنه \emph{ذو~$\OLId{latin-ordinary}{n}$ من
  المواضع}. والرابط ذو الرتبة~$\OLMathNumeral{0}$ يسمى أيضًا \emph{ثابتًا}؛ وذو
  الرتبة~$\OLMathNumeral{1}$ يسمى \emph{أحاديًا}، وذو الرتبة~$\OLMathNumeral{2}$ \emph{ثنائيًا}.
\end{defn}

\begin{ex}
  اللغة القياسية لمنطق القضايا~$\Lang L_\OLMathNumeral{0}$ مؤلفة من الروابط الآتية،
  ومع كل منها رتبته:
  $\lfalse$~($\OLMathNumeral{0}$)،
  $\lnot$~($\OLMathNumeral{1}$)،
  $\land$~($\OLMathNumeral{2}$)،
  $\lor$~($\OLMathNumeral{2}$)،
  $\lif$~($\OLMathNumeral{2}$). وأكثر المنطقات الآتية مبني على هذه اللغة.
  ونرمز بـ~$\Lang L_\OLMathNumeral{0}^{-}$ إلى جزئها الخالي من الثوابت، المتولد من
  متغيرات القضايا والروابط~$\lnot$ و$\land$ و$\lor$ و$\lif$ وحدها.
  % تصحيح مصدر مشترك (MVL-L0-I1): عُرّف الجزء الخالي من الثوابت لئلا تنسب مصفوفة إلى L_0 من غير تفسير الثابت lfalse.
  وقد يجري الاصطلاح في بعض المنطقات برمز آخر لرابط مألوف؛ ففي منطق
  الضرب، مثلًا، يرمز إلى الاقتران غالبًا بـ~$\odot$ لا بـ~$\land$.
  وقد تدعو الحاجة إلى مؤثر زائد، كمؤثر التعيّن~$\triangle$، وهو ذو
  موضع واحد.
\end{ex}

\end{document}
